\KOMAoption{paper}{a4, portrait}
\areaset{180mm}{267mm}
\fancyheadoffset{0pt} % обновляем хедер и футер

\newgeometry{
  includehead,
  includefoot,
  nomarginpar,
  top=6mm,
  bottom=6mm,
  left=22mm,
  right=10mm,
  headheight=14.5pt,
  headsep=3mm,        % Добавить headsep
  footskip=7mm,
}
%\headsep=0.3cm
%\setlength{\footskip}{30pt}
%\topmargin=-2.6cm % Команда \topmargin задаёт верхнее поле страницы. При этом поле отсчитывается не от левого края листа, а от линии, параллельной краю листа и отстоящей от него на 1 дюйм
%\oddsidemargin=-0.3cm % отступ слева https://zns.susu.ru/IT/latex/Marking/marking.html

%\footskip=0cm
%\setlength{\headheight}{1cm}
%\setlength{\footheight}{1cm}
%\setlength{\textheight}{205mm}

\color{uniblue}{\section[(рекомендуемое) Элементы функциональных схем]{(рекомендуемое)\\Элементы функциональных схем}\label{app:func_desc}}
\color{black}


%\small % сделаем шрифт в таблице поменьше

\begin{longtable}{
|>{\centering\arraybackslash}m{5.5cm}
|>{\centering\arraybackslash}m{5.5cm}
|>{\centering\arraybackslash}m{5.5cm}
|}

%\caption{Принятые обозначения элементов функциональных схем\vspace{-0.5\baselineskip}}\label{appEl:tbl1}\\ % выравнивает надпись посередине таблицы
\caption{Принятые обозначения элементов функциональных схем\hfill\vspace{-0.5\baselineskip}}\label{appEl:tbl1}\\ 

\hline
\rowcolor{gray!30}
Графическое отображение на схеме & Назначение & Примечание \\ 
\hline
\endfirsthead
%\caption{\emph{Продолжение\vspace{-0.5\baselineskip}}} \\ % Отступ обозначения таблицы от таблицы и выравнивание надписи слева
\caption*{\hspace{3pt}\emph{Продолжение таблицы \ref{appEl:tbl1}\hfill\vspace{-0.5\baselineskip}}} \\ % сделано по ГОСТ 2.105 п.6.8.7
\hline
\rowcolor{gray!30}
Графическое отображение на схеме & Назначение & Примечание \\ 
\hline
\endhead

\hline
\endfoot

\hline
\endlastfoot

\centering {
\includegraphics[width=0.2\textwidth,height=0.2\textheight,keepaspectratio]{\genpath/Элементы ФС/_latex/img/1_analog_canal.pdf}}& \centering Аналоговый канал & \centering \arraybackslash\\
\hline

\centering {\includegraphics[width=0.2\textwidth,height=0.2\textheight,keepaspectratio]{\genpath/Элементы ФС/_latex/img/2_po.pdf}} & \centering Пусковой (измерительный орган) & \centering \arraybackslash   \\
\hline

\centering {\includegraphics[width=0.2\textwidth,height=0.2\textheight,keepaspectratio]{\genpath/Элементы ФС/_latex/img/3_out_sig.pdf}} & \centering Внешний входной сигнал (из матрицы входов/выходов) & \centering \arraybackslash   \\
\hline

\centering {\includegraphics[width=0.2\textwidth,height=0.2\textheight,keepaspectratio]{\genpath/Элементы ФС/_latex/img/4_inner.pdf}} & \centering Входной сигнал от внутренней логики & \centering \arraybackslash   \\
\hline

\centering {\includegraphics[width=0.2\textwidth,height=0.2\textheight,keepaspectratio]{\genpath/Элементы ФС/_latex/img/5_out_inner.pdf}} & \centering Выходной сигнал внутренней логики & \centering \arraybackslash   \\
\hline

\centering {\includegraphics[width=0.2\textwidth,height=0.2\textheight,keepaspectratio]{\genpath/Элементы ФС/_latex/img/6_virt_key.pdf}} & \centering Входной сигнал от <<Виртуального ключа>> & \centering \arraybackslash   \\
\hline

\centering {\includegraphics[width=0.2\textwidth,height=0.2\textheight,keepaspectratio]{\genpath/Элементы ФС/_latex/img/7_virt_but.pdf}} & \centering Входной сигнал от <<Виртуальной кнопки>> & \centering \arraybackslash   \\
\hline

\centering {
\includegraphics[width=0.2\textwidth,height=0.2\textheight,keepaspectratio]{\genpath/Элементы ФС/_latex/img/8_switch.pdf}} & \centering Программный переключатель & \centering \arraybackslash   \\
\hline

\centering {\includegraphics[width=0.2\textwidth,height=0.2\textheight,keepaspectratio]{\genpath/Элементы ФС/_latex/img/9_1_and.pdf}} & \centering Логический элемент <<И>> & \centering \arraybackslash 
  {\includegraphics[width=0.1\textwidth,height=0.1\textheight,keepaspectratio]{\genpath/Элементы ФС/_latex/img/9_2_and.pdf}} \\
\hline

\centering {\includegraphics[width=0.2\textwidth,height=0.2\textheight,keepaspectratio]{\genpath/Элементы ФС/_latex/img/10_1_or.pdf}} & \centering Логический элемент <<ИЛИ>> & \centering \arraybackslash 
  {\includegraphics[width=0.1\textwidth,height=0.1\textheight,keepaspectratio]{\genpath/Элементы ФС/_latex/img/10_2_or.pdf}} \\
\hline

\centering {\includegraphics[width=0.2\textwidth,height=0.2\textheight,keepaspectratio]{\genpath/Элементы ФС/_latex/img/11_1_xor.pdf}} & \centering Логический элемент <<ИСКЛЮЧАЮЩЕЕ ИЛИ>> & \centering \arraybackslash 
  {\includegraphics[width=0.1\textwidth,height=0.1\textheight,keepaspectratio]{\genpath/Элементы ФС/_latex/img/11_2_xor.pdf}} \\
\hline

\centering {\includegraphics[width=0.2\textwidth,height=0.2\textheight,keepaspectratio]{\genpath/Элементы ФС/_latex/img/12_1_sr_mem.pdf}} & \centering SR-триггер & \multirow{2}{*} 
  {\includegraphics[width=0.1\textwidth,height=0.1\textheight,keepaspectratio]{\genpath/Элементы ФС/_latex/img/12_2.pdf}} \\
\cline{1-2}

\centering {\includegraphics[width=0.2\textwidth,height=0.2\textheight,keepaspectratio]{\genpath/Элементы ФС/_latex/img/13_1_sr.pdf}} & \centering SR-триггер с сохранением состояния после исчезновения питания & \centering \arraybackslash   \\
\hline

\centering {\includegraphics[width=0.2\textwidth,height=0.2\textheight,keepaspectratio]{\genpath/Элементы ФС/_latex/img/14_1_rs_mem.pdf}} & \centering RS-триггер & \multirow{2}{*} 
  {\includegraphics[width=0.1\textwidth,height=0.1\textheight,keepaspectratio]{\genpath/Элементы ФС/_latex/img/14_2.pdf}} \\
\cline{1-2}

\centering {\includegraphics[width=0.2\textwidth,height=0.2\textheight,keepaspectratio]{\genpath/Элементы ФС/_latex/img/15_1_rs.pdf}} & \centering RS-триггер с сохранением состояния после исчезновения питания & \centering \arraybackslash   \\
\hline

\centering {\includegraphics[width=0.2\textwidth,height=0.2\textheight,keepaspectratio]{\genpath/Элементы ФС/_latex/img/16_not.pdf}}& \centering Логический элемент <<НЕ>> & \centering \arraybackslash\\
\hline

\centering {\includegraphics[width=0.2\textwidth,height=0.2\textheight,keepaspectratio]{\genpath/Элементы ФС/_latex/img/17_1_ton.pdf}} & \centering Выдержка времени на срабатывание (регулируемая) & \multirow{2}{*} 
  {\includegraphics[width=0.23\textwidth,height=0.23\textheight,keepaspectratio]{\genpath/Элементы ФС/_latex/img/17_2_ton.pdf}} \\
\cline{1-2}

\centering {\includegraphics[width=0.2\textwidth,height=0.2\textheight,keepaspectratio]{\genpath/Элементы ФС/_latex/img/18.pdf}} & \centering Выдержка времени на срабатывание (нерегулируемая) & \centering \arraybackslash   \\
\hline

\centering {\includegraphics[width=0.2\textwidth,height=0.2\textheight,keepaspectratio]{\genpath/Элементы ФС/_latex/img/19_1_tof.pdf}} & \centering Выдержка времени на возврат (регулируемая) & \multirow{2}{*} 
  {\includegraphics[width=0.23\textwidth,height=0.23\textheight,keepaspectratio]{\genpath/Элементы ФС/_latex/img/20.pdf}} \\
\cline{1-2}

\centering {\includegraphics[width=0.2\textwidth,height=0.2\textheight,keepaspectratio]{\genpath/Элементы ФС/_latex/img/19_2_tof.pdf}} & \centering Выдержка времени на возврат (нерегулируемая) & \centering \arraybackslash   \\
\hline

\centering {\includegraphics[width=0.2\textwidth,height=0.2\textheight,keepaspectratio]{\genpath/Элементы ФС/_latex/img/21_1_pls.pdf}} & \centering Одновибратор по переднему фронту без перезапуска (регулируемый) & \multirow{2}{*} 
  {\includegraphics[width=0.23\textwidth,height=0.23\textheight,keepaspectratio]{\genpath/Элементы ФС/_latex/img/21_2_pls.pdf}} \\
\cline{1-2}

\centering {\includegraphics[width=0.2\textwidth,height=0.2\textheight,keepaspectratio]{\genpath/Элементы ФС/_latex/img/22.pdf}} & \centering Одновибратор по переднему фронту без перезапуска (нерегулируемый) & \centering \arraybackslash   \\
\hline

\centering {\includegraphics[width=0.2\textwidth,height=0.2\textheight,keepaspectratio]{\genpath/Элементы ФС/_latex/img/23_1.pdf}} & \centering Одновибратор по переднему фронту со сбросом (регулируемый) & \centering \arraybackslash 
  {\includegraphics[width=0.2\textwidth,height=0.2\textheight,keepaspectratio]{\genpath/Элементы ФС/_latex/img/23_2.pdf}} \\
\hline

\centering {\includegraphics[width=0.2\textwidth,height=0.2\textheight,keepaspectratio]{\genpath/Элементы ФС/_latex/img/24_1.pdf}} & \centering Одновибратор по заднему фронту без перезапуска (нерегулируемый) & \centering \arraybackslash 
  {\includegraphics[width=0.2\textwidth,height=0.2\textheight,keepaspectratio]{\genpath/Элементы ФС/_latex/img/24_2.pdf}} \\
\hline

\centering {\includegraphics[width=0.2\textwidth,height=0.2\textheight,keepaspectratio]{\genpath/Элементы ФС/_latex/img/25.pdf}} & \centering Выбор максимального значения & \centering \arraybackslash   \\
\hline

\centering {\includegraphics[width=0.2\textwidth,height=0.2\textheight,keepaspectratio]{\genpath/Элементы ФС/_latex/img/26.pdf}} & \centering Выбор минимального значения & \centering \arraybackslash   \\
\hline

\centering {\includegraphics[width=0.2\textwidth,height=0.2\textheight,keepaspectratio]{\genpath/Элементы ФС/_latex/img/27.pdf}} & \centering Пороговый элемент с гистерезисом (сравнение с уставкой) & \centering \arraybackslash   \\
\hline

\end{longtable} 

%clearpage 
\restoregeometry